\clearpage\begingroup\SourceLeaf{642b}
\ArithmeticHead{7}{R}{I. NACHWEIS DER ERDROTATION.}
\makeatletter\tagsleft@true\let\veqno\@@leqno\makeatother
\renewcommand{\thefootnote}{\arabic{footnote})}\setcounter{footnote}{7}
\begin{center}\textls[100]{2. Die \textsc{Gauss}sche Arbeit.}\end{center}
\noindent\hspace*{\parindent}Die \textsc{Gauss}sche Behandlungsweise unterscheidet sich von der \textsc{Laplace}schen
insbesondere durch die geschickte Wahl der Bezugssysteme und die elegantere,
suggestivere Durchführung; \textsc{Gauss} selbst scheint dabei die \textsc{Laplace}sche Abhandlung
nicht gekannt zu haben. \textsc{Gauss} führt zwei Koordinatensysteme ein:
1. ein festes System $X,Y,Z$, bei dem die Ebene $Z=0$ parallel zum Äquator
durch den Ausgangspunkt des fallenden Körpers geht, während die Ebene
$Y=0$ mit der Meridianebene dieses Punktes übereinstimmt, und der Anfangs\nobreakdash-
\clearpage\SourceLeaf{643}\ArithmeticHead{8}{L}{H. GEPPERT, ÜBER GAUSS’ ARBEITEN ZUR MECHANIK UND POTENTIALTHEORIE.}
\noindent punkt $O$ der entsprechende Punkt der Erdachse ist; 2. ein mit der Erde mitrotierendes
System $x,y,z$, dessen Anfangspunkt mit dem der Fallbewegung
übereinstimmt, und bei dem die Ebene $y=0$ senkrecht auf der Richtung der
scheinbaren Schwere steht, sodass die drei Achsen bezw. nach Süden, Osten
und dem scheinbaren Zenith weisen. Die Formeln für die Transformation
des einen Systems in das andere sind leicht zu entwickeln; in sie geht naturgemäss
die Zeit $t$ des Falles (und damit der Erdrotation) ein. \textsc{Gauss} entwickelt
sofort die Bewegungsgleichungen des fallenden Körpers bezüglich des festen
Systems unter der Annahme eines zum Geschwindigkeitsquadrat proportionalen
Luftwiderstandes und gewinnt daraus die Gleichungen im bewegten System.
Letztere zeigen nun, dass auf den ruhenden Körper nur die von der Zentrifugalkraft
affizierte Schwerkraft wirkt, die die bekannte Lotabweichung zur
Folge hat, dass hingegen für den bewegten Körper noch eine Kraft (\textsc{Coriolis}kraft)
hinzutritt, die allein es ermöglicht, die Rotation der Erde sichtbar
zu machen. \textsc{Gauss} ist somit wohl der erste, der dieser zusätzlichen \textsc{Coriolis}kraft
(die erst viel später in die allgemeine Theorie einging) die fügliche Beachtung
schenkte. Unter dem Vorbehalt eines im Verhältnis zum Erdradius
kleinen Fallraumes, wie er bei Versuchen tatsächlich vorliegt, ist die scheinbare
Schwere $g$ während der Bewegung der Grösse und Richtung nach konstant,
und daher gewinnen die Bewegungsgleichungen die einfache Form
\[
\begin{aligned}
\ddot x-2n\sin\varphi\,\dot y+Mu\dot x&=0,\\
\ddot y+2n\sin\varphi\,\dot x+2n\cos\varphi\,\dot z+Mu\dot y&=0,\\
\ddot z-2n\cos\varphi\,\dot y+Mu\dot z+g&=0,
\end{aligned}\tag{7}
\]
\noindent worin $\varphi$ die geographische Breite des Beobachtungsortes, $n=\frac{1}{13713}$ die Winkelgeschwindigkeit
der Erdrotation, $M$ einen konstanten Reibungsfaktor und
$u=\sqrt{\dot x^{2}+\dot y^{2}+\dot z^{2}}$ die totale Relativgeschwindigkeit des Körpers bedeuten. Diese
linearen Differentialgleichungen sind unter den vorgegebenen Anfangsbedingungen
leicht zu integrieren; entwickelt man die Lösungen nach Potenzen
der kleinen Grössen $n$ und $M$, so finden sich schliesslich die folgenden Ausdrücke:
\clearpage\SourceLeaf{644}\ArithmeticHead{9}{R}{I. NACHWEIS DER ERDROTATION.}
\HistoricalFoot{X 2 Abh. 7.}{2}
\[
\left\{\begin{aligned}
&x=\text{südl. Abweichung vom scheinbaren Lote}\\[-2pt]
&\hspace{2em}=\tfrac16 g\sin\varphi\cos\varphi\mathpunct{.}\,n^{2}t^{4},\\
&y=\text{östl. Abweichung vom scheinbaren Lote}\\[-2pt]
&\hspace{2em}=\tfrac13 g\cos\varphi\mathpunct{.}\,nt^{3}-\tfrac1{12}g^{2}\cos\varphi\mathpunct{.}\,Mnt^{5},\\
&z=\text{Fallstrecke in Richtung des scheinbaren Lotes}\\[-2pt]
&\hspace{2em}=\tfrac12 gt^{2}-\tfrac1{12}Mg^{2}t^{4}.
\end{aligned}\right.\tag{8}
\]
\noindent Man sieht, dass der Wert von $y$ mit der \textsc{Laplace}schen Formel (4) (worin $\theta$
mit $\frac\pi2-\varphi$, $m$ mit $M$ zu identifizieren ist) übereinstimmt; auch \textsc{Gauss} findet,
wie \textsc{Laplace}, eine völlig unterhalb der Beobachtungsgrenze liegende südliche
Abweichung.

\textsc{Gauss} selbst machte sogleich die Anwendung seiner Theorie auf die
\textsc{Benzenberg}schen Versuche und fand den oben angegebenen theoretischen Wert
der östlichen Abweichung von $8{,}7$ mm. Den besprochenen Aufsatz hat \textsc{Gauss}
im August 1803 auf Wunsch von \textsc{Benzenberg} neu durchgearbeitet, damit er
in dessen Buche veröffentlicht würde. Während des ganzen Monats März 1803
dreht sich die Korrespondenz zwischen \textsc{Gauss} und \textsc{Olbers} um die Entscheidung,
welche ihrer beiden Falltheorien die richtige sei. Die \textsc{Olbers}schen
Überlegungen, die zum grossen Teil mit den schon von \textsc{Guglielmini} in seinem
Buche angestellten Betrachtungen übereinstimmen, sind wegen ihrer elementaren
Natur besonders interessant. \textsc{Olbers} machte jedoch zwei Annahmen,
die sich beide als irrig erweisen: erstens, dass die Schwererichtungen in allen
Punkten der Fallbewegung im absoluten Raume parallel seien, und zweitens,
dass der fallende Körper in der Ebene bleibe, die durch Ausgangspunkt und
Erdmittelpunkt geht und auf dem anfänglichen Meridian senkrecht steht. Die
erste Annahme hat zur Folge, dass die östliche Abweichung sich als das
Anderthalbfache ihres wahren Wertes errechnet, die zweite, zunächst sehr naheliegende
Annahme widerlegt sich aus den \textsc{Gauss}schen Grundgleichungen (7);
denn infolge der Reibung wird der fallende Körper aus der genannten Ebene
mit der Kraft $Muny\sin\varphi$ nach Norden abgedrängt, wodurch gerade die von
\textsc{Olbers} errechnete südliche Abweichung kompensiert wird. Nachdem \textsc{Gauss}
diese Fehler \textsc{Olbers} einzeln nachgewiesen und ihn überzeugt hatte, gab \textsc{Olbers}
eine berichtigte Zusammenfassung seiner elementaren Theorie in \textsc{Benzenbergs}
Buch (S. 372—383); auch \textsc{Guglielmini} hat seine ursprüngliche Theorie nach
den \textsc{Laplace}schen Angaben berichtigt\footnote{Vergl. \textsc{Guglielminis} Brief an \textsc{Benzenberg} vom 23. März 1803; \textsc{Benzenberg} a.\,a.\,O., S. 384—387.}.
\clearpage\endgroup
